% Part: first-order-logic
% Chapter: syntax-and-semantics
% Section: main-operator

\documentclass[../../../include/open-logic-section]{subfiles}

\begin{document}

\olfileid{fol}{syn}{mai}

\olsection{\printtoken{S}{main operator} sawijining Formula}

\begin{explain}
Kerep migunani kanggo ngrembug operator pungkasan kang digunakake
nalika mbangun !!a{formula}~$!A$. Operator iki diarani \emph{operator
  utama} saka~$!A$. Sacara intuitif, iki operator kang paling njaba
ing $!A$. Umpamane, operator utama saka $\lnot !A$ yaiku $\lnot$,
operator utama saka $(!A \lor !B)$ yaiku $\lor$, lan sateruse.
\end{explain}


\begin{defn}[!!^{main operator}]
\ollabel{def:main-op}
\emph{!!{main operator}} saka !!a{formula}~$!A$
ditetepake kaya mangkene:
\begin{enumerate}
\item \indcase*{!A}{!A}{$\indfrm$ ora nduweni !!{main operator}.}

\tagitem{prvNot}{\indcase{!A}{\lnot !B}{!!{main operator} saka $\indfrm$
  yaiku~$\lnot$.}}{}

\tagitem{prvAnd}{\indcase{!A}{(!B \land !C)}{!!{main operator} saka
  $\indfrm$ yaiku~$\land$.}}{}

\tagitem{prvOr}{\indcase{!A}{(!B \lor !C)}{!!{main operator} saka
  $\indfrm$ yaiku~$\lor$.}}{}

\tagitem{prvIf}{\indcase{!A}{(!B \lif !C)}{!!{main operator} saka
  $\indfrm$ yaiku~$\lif$.}}{}

\tagitem{prvIff}{\indcase{!A}{(!B \liff !C)}{!!{main operator} saka
  $\indfrm$ yaiku~$\liff$.}}{}

\tagitem{prvAll}{\indcase{!A}{\lforall[x][!B]}{!!{main operator}
  saka $\indfrm$ yaiku~$\lforall$.}}{}

\tagitem{prvEx}{\indcase{!A}{\lexists[x][!B]}{!!{main operator} saka
  $\indfrm$ yaiku~$\lexists$.}}{}
\end{enumerate}
\end{defn}

Ing saben kasus, kang kita maksud yaiku \emph{muncule} tartamtu
saka !!{main operator} kang dituduhake ing formula. Umpamane, amarga
formula $((!D \lif !E) \lif (!E \lif !D))$ awujud $(!B \lif !C)$
kanthi $!B$ yaiku $(!D \lif !E)$ lan $!C$ yaiku $(!E \lif !D)$,
muncule $\lif$ kang kapindho iku !!{main operator}.

\begin{explain}
Iki \emph{definisi rekursif} sawijining fungsi kang nggayutake kabeh
!!{formula}s nonatomik karo muncule !!{main operator}. Amarga
!!{formula}s ditetepake kanthi induktif, saben
!!{formula}~$!A$ nyukupi salah siji kasus ing \olref{def:main-op}.
Iki njamin yen kanggo saben !!{formula} nonatomik~$!A$ ana !!a{main
  operator}. Amarga saben !!{formula} mung nyukupi siji saka
syarat-syarat iki, lan amarga !!{formula}s kang luwih cilik kang
digunakake kanggo mbangun $!A$ ditemtokake kanthi siji cara ing
saben kasus, muncule !!{main
  operator} saka~$!A$ iku tunggal, mula kita wis netepake sawijining
fungsi.
\end{explain}

Kita menehi jeneng marang !!{formula}s miturut \olref{tab:main-op}
gumantung marang simbol kang dadi !!{main operator}.
\iftag{defNot,defOr,defAnd,defIf,defIff,defTrue,defFalse,defEx,defAll}
{ Nanging, elinga yen operator kang ditetepake ora resmi muncul ing
!!{formula}s. Operator kasebut mung cekakan, mula sacara resmi ora
bisa dadi operator utama sawijining formula. Mula, ing bukti ngenani
kabeh !!{formula}s, operator kasebut ora perlu ditangani dhewe.}

\begin{table}[!h]
\centering
\begin{tabular}{c | c | c}
!!^{main operator} & Jinis !!{formula} & Tuladha\\
\hline
ora ana & atomik (!!{formula}) &
\iftag{prvFalse}{$\lfalse$,}{}
\iftag{prvTrue}{$\ltrue$,}{}
$\Atom{R}{t_1, \dots, t_n}$,
$\eq[t_1][t_2]$\\
$\lnot$ & negasi & $\lnot !A$ \\
$\land$ & konjungsi & $(!A \land !B)$ \\
$\lor$ & disjungsi & $(!A \lor !B)$ \\
$\lif$ & !!{conditional} & $(!A \lif !B)$ \\
$\liff$ & !!{biconditional} & $(!A \liff !B)$ \\
$\lforall[][]$ & universal (!!{formula})& $\lforall[x][!A]$ \\
$\lexists[][]$ & eksistensial (!!{formula})& $\lexists[x][!A]$
\end{tabular}
\caption{Operator utama lan jeneng !!{formula}s. Cathetan koreksi sumber
OLPL-066: ing telung larik konjungsi, disjungsi, lan implikasi, kurung
tutup dipindhah mlebu wates rumus supaya pasangan kurunge pepak.}
\ollabel{tab:main-op}
\end{table}

\end{document}
